%% kapa: exam=1 round=28 year=2017 period=1 form=A subject=civil_law subject_ko=민법 seq=3 number=3
%% source: Q-Net 감정평가사 기출문제 (cst00309) artlSeq=5199603 file=extracted/1008/5199603/2017년 제28회 감정평가사 1차시험 문제지/1교시 A형.pdf
%% conversion: kordoc 4.24.51 -> tools/gen1.py (auto); review=proofread-sonnet; flags=none
\begin{question}{3}
제한능력자에 관한 설명으로 옳은 것만을 모두 고른 것은? (다툼이 있으면 판례에 따름)
\begin{framedtext}
ㄱ. 만18세의 미성년자가 자기의 월 근로소득 범위 내에서 신용구매계약을 체결한 경우, 그 신용구매계약은 처분허락을 받은 재산범위 내의 처분행위에 해당한다. \par
ㄴ. 한정후견인의 동의가 필요한 법률행위를 피한정후견인이 한정후견인의 동의 없이 하였을 때에는 그것이 일상생활에 필요하고 그 대가가 과도하지 아니한 법률행위가 아닌 경우 그 법률행위를 취소할 수 있다. \par
ㄷ. 제한능력자가 아직 능력자가 되지 못한 경우에도 그 상대방은 그에게 1개월 이상의 기간을 정하여 추인 여부의 확답을 촉구할 수 있다. \par
ㄹ. 제한능력자와 계약을 맺은 선의의 상대방은 추인이 있기 전까지 의사표시를 거절할 수 있다.
\end{framedtext}
\begin{choices}
  \item ㄱ, ㄴ
  \item ㄱ, ㄷ
  \item ㄴ, ㄷ
  \item ㄴ, ㄹ
  \item ㄷ, ㄹ
\end{choices}
\end{question}
\answer{1}
